\begin{figure}[H]
    \centering
    \begin{tikzpicture}
        \begin{loglogaxis}[
            width=0.88\textwidth, height=7.2cm,
            xlabel={$n$}, ylabel={tempo mediano (ms)},
            legend pos=north west, legend cell align=left,
            legend style={font=\scriptsize}, grid=both, font=\footnotesize,
            log basis x=10, log basis y=10]
            \addplot+[mark=*] table[x=n,y=t,col sep=comma] {results/plot/mcm-ek.csv};
            \addlegendentry{Edmonds-Karp $O(V E^2)$}
            \addplot+[mark=square*] table[x=n,y=t,col sep=comma] {results/plot/mcm-aug.csv};
            \addlegendentry{Caminhos aumentantes $O(VE)$}
            \addplot+[mark=triangle*] table[x=n,y=t,col sep=comma] {results/plot/mcm-ks.csv};
            \addlegendentry{Karp-Sipser (heur.)}
            \addplot+[mark=diamond*] table[x=n,y=t,col sep=comma] {results/plot/mcm-hk.csv};
            \addlegendentry{Hopcroft-Karp $O(E\sqrt{V})$}
            \addplot+[mark=o] table[x=n,y=t,col sep=comma] {results/plot/mcm-greedy.csv};
            \addlegendentry{Gulosa (heur.)}
        \end{loglogaxis}
    \end{tikzpicture}
    \caption{Tempo de execução em escala log-log para o grafo denso (cardinalidade
    máxima bipartida). A inclinação de cada reta corresponde ao expoente empírico
    de crescimento; a de Edmonds-Karp, próxima de 3, destaca-se das demais.}
    \label{fig:cresc-mcm}
\end{figure}
